\section{Lesson 4: Crossing State Lines Produces Compact but Fragmented Districts}
\label{sec:lesson4}

National redistricting without state lines (Paper E.3) is the purest
implementation of geometric optimization: treat the continental United States
as a single graph of approximately 73,000 census tracts and partition it into
435 compact contiguous districts using the METIS minimum-cut algorithm, without
any constraint that district boundaries respect state lines
\citep{dellaLibera2026e3}.  The design lesson from this experiment is a
precise quantification of the federalism compactness cost: the state-boundary
constraint sacrifices approximately 12\% of achievable mean Polsby--Popper
compactness in exchange for state-based representation.

\subsection{The Compactness Gain}

The 50-state baseline DIA achieves a mean Polsby--Popper score of 0.367.
The national redistricting experiment achieves 0.411, an improvement of
$\Delta\mathrm{PP} = +0.044$, or 12.0\% of the baseline value
\citep{dellaLibera2026e3}.

This improvement is not uniform across the country.  It is concentrated in
states where state boundaries happen to bisect naturally compact geographic
units: the Appalachian ridge that splits the Virginia--West Virginia--North
Carolina border region; the Mississippi River basin that straddles multiple
state lines in the upper Midwest; the suburban sprawl of metropolitan areas
that cross state boundaries (Kansas City, Philadelphia, Cincinnati, Washington
DC).  In these areas, removing the state constraint allows the algorithm to
draw districts that follow natural geographic contours rather than arbitrary
political lines.

Approximately 60\% of districts are unchanged or minimally changed by removing
the state constraint: they would fall entirely within one state under any
reasonable compact-district drawing.  The compactness gain is concentrated in
the remaining 40\%, which gain an average of $\Delta\mathrm{PP} = +0.110$
--- a substantial improvement in the districts where the state constraint
was binding.

\subsection{The County Preservation Cost}

The compactness improvement comes at a significant cost in county preservation.
The baseline DIA preserves 38\% of counties whole (no tract splits across
district lines).  The national redistricting experiment preserves only 23\%
of counties whole --- a 40\% reduction in county-preservation rate.

This is counterintuitive at first: if the algorithm is optimizing for
compactness without state-line constraints, why does county preservation
decline?  The answer is that state boundaries often happen to coincide with
county boundaries --- state lines commonly run along county borders, especially
in the South and Midwest where counties were established contemporaneously with
state boundaries.  When the algorithm is free to cross state lines, it draws
districts that cross county lines in new places, fragmenting counties that
had been preserved under the state-constrained baseline.

The tradeoff between compactness and county preservation is thus partially
mediated by the state-boundary constraint: state boundaries impose compactness
costs in some places while incidentally protecting county integrity in others.
Removing the state constraint improves geometric compactness while degrading
county coherence.

\subsection{Political Fragmentation}

Beyond county preservation, national redistricting produces a specific type of
political fragmentation that does not arise in the state-constrained baseline:
districts that cross state lines represent constituents governed by different
state laws.

A district spanning the Kansas--Missouri border would represent residents of
two states with different income tax rates, different Medicaid eligibility
thresholds, different state legislative districts, different US Senate
representation, and different electoral calendars.  The House member from such
a district would face constituent questions --- about state programs, state
regulations, state courts --- that do not map cleanly onto her congressional
role.  She would also face a reelection challenge defined by two different
state primary calendars and two different state party structures.

Paper E.3 identifies 87 cross-state districts in the national map --- districts
that draw at least 10\% of their population from a different state than the
majority of their population \citep{dellaLibera2026e3}.  These 87 districts
represent a substantial fraction of the 435 total (20\%) and are concentrated
in the densely settled Northeast and the suburban Midwest.

\subsection{The Federalism Cost: A Precise Measurement}

The central contribution of E.3 to the design lessons is that it makes the
federalism cost of redistricting legible.  Before this experiment, the cost of
respecting state boundaries in redistricting was unmeasured; the conventional
wisdom was simply that state-based redistricting is constitutionally required
and therefore not worth analyzing as a design choice.

The finding that the federalism constraint costs approximately 12\% mean
compactness is both modest enough to be acceptable and substantial enough to
be worth knowing.  It is modest in the sense that the 0.367 baseline is
already substantially higher than the compactness of current commission-drawn
maps (which average approximately 0.28--0.31 in most states); the 12\%
improvement would raise the algorithmic baseline to 0.411, an improvement at
the margin.  It is substantial in the sense that 12\% of mean PP corresponds
to visibly more regular district shapes in the cross-state border regions where
the improvement is concentrated.

The implication for reform advocates is that the compactness argument for
national redistricting is relatively weak: a 12\% improvement is not
transformative.  The proportionality argument --- that national redistricting
might reduce the sorting-driven seat gap by drawing districts that cross
state-level sorting concentrations --- is potentially stronger but is not
supported by E.3's empirical results (see Lesson~1 above).
